\documentclass[11pt]{article}

\usepackage[a4paper,margin=27mm]{geometry}
\usepackage{amsmath,amssymb,amsthm,mathtools}
\usepackage{microtype}
\usepackage[hidelinks]{hyperref}

\newtheorem{theorem}{Theorem}
\newtheorem{lemma}{Lemma}
\theoremstyle{remark}
\newtheorem*{remark}{Remark}

\newcommand{\Z}{\mathbb Z}
\newcommand{\N}{\mathbb N}

\title{An Infinite Pell-Type Family for\texorpdfstring{\\[2mm]
$y^2+x^2y+xz^2+y-1=0$}{ y2+x2y+xz2+y-1=0}}
\author{}
\date{}

\begin{document}
\maketitle
\vspace{-2.5em}

\begin{abstract}
We give a completely explicit infinite family of integer solutions to
\[
 y^2+x^2y+xz^2+y-1=0.
\]
All solutions in the family have the fixed coordinate $x=-55$.  The proof is an elementary Pell-type norm argument: one seed point is propagated by an integral linear transformation preserving the quadratic form $U^2-55V^2$.
\end{abstract}

\section{The infinite family}

Define integer sequences $(U_n)_{n\geq 0}$ and $(V_n)_{n\geq 0}$ by
\begin{equation}\label{eq:recurrence}
 U_0=1575,\qquad V_0=59,
\end{equation}
and
\begin{equation}\label{eq:recurrence-step}
 \begin{pmatrix}U_{n+1}\\ V_{n+1}\end{pmatrix}
 =
 \begin{pmatrix}89&660\\12&89\end{pmatrix}
 \begin{pmatrix}U_n\\V_n\end{pmatrix};
\end{equation}
that is,
\[
 U_{n+1}=89U_n+660V_n,
 \qquad
 V_{n+1}=12U_n+89V_n.
\]
Equivalently,
\begin{equation}\label{eq:closed-form}
 U_n+V_n\sqrt{55}
 =(1575+59\sqrt{55})(89+12\sqrt{55})^n.
\end{equation}

\begin{lemma}[Norm preservation]\label{lem:norm}
For all integers $U,V$, if
\[
 U'=89U+660V,
 \qquad
 V'=12U+89V,
\]
then
\[
 (U')^2-55(V')^2=U^2-55V^2.
\]
\end{lemma}

\begin{proof}
Direct expansion gives
\begin{align*}
 (U')^2-55(V')^2
 &=(89U+660V)^2-55(12U+89V)^2\\
 &=(89^2-55\cdot12^2)U^2
   +2\cdot89(660-55\cdot12)UV\\
 &\qquad +(660^2-55\cdot89^2)V^2.
\end{align*}
The three coefficients are respectively
\[
 89^2-55\cdot12^2=1,
 \qquad
 2\cdot89(660-55\cdot12)=0,
 \qquad
 660^2-55\cdot89^2=-55.
\]
Hence the displayed expression is $U^2-55V^2$.
\end{proof}

\begin{theorem}\label{thm:main}
For every $n\in\Z_{\geq0}$ and every $\varepsilon,\delta\in\{\pm1\}$,
\begin{equation}\label{eq:family}
 \boxed{
 (x,y,z)=\bigl(-55,\,\varepsilon U_n-1513,\,\delta V_n\bigr)
 }
\end{equation}
is an integer solution of
\begin{equation}\label{eq:original}
 y^2+x^2y+xz^2+y-1=0.
\end{equation}
Moreover, the triples with $\varepsilon=\delta=1$ are pairwise distinct.  Consequently, \eqref{eq:original} has infinitely many integer solutions.
\end{theorem}

\begin{proof}
First,
\begin{equation}\label{eq:seed-norm}
 1575^2-55\cdot59^2
 =2\,480\,625-191\,455
 =2\,289\,170
 =1513^2+1.
\end{equation}
Lemma~\ref{lem:norm}, applied inductively to \eqref{eq:recurrence-step}, therefore yields
\begin{equation}\label{eq:all-norms}
 U_n^2-55V_n^2=1513^2+1
 \qquad(n\geq0).
\end{equation}

Now fix $n\geq0$ and $\varepsilon,\delta\in\{\pm1\}$, and set
\[
 x=-55,
 \qquad
 y=\varepsilon U_n-1513,
 \qquad
 z=\delta V_n.
\]
Because $x^2+1=3026=2\cdot1513$, the left-hand side of \eqref{eq:original} is
\begin{align*}
 y^2+x^2y+xz^2+y-1
 &=y^2+3026y-55z^2-1\\
 &=(y+1513)^2-55z^2-(1513^2+1)\\
 &=U_n^2-55V_n^2-(1513^2+1)=0
\end{align*}
by \eqref{eq:all-norms}.  Thus every triple in \eqref{eq:family} is an integer solution.

It remains only to prove that infinitely many of these triples are distinct.  Since $U_0,V_0>0$ and every entry of the recurrence matrix in \eqref{eq:recurrence-step} is positive, induction gives $U_n,V_n>0$ for all $n$.  Hence
\[
 V_{n+1}=12U_n+89V_n>V_n.
\]
Thus $(V_n)_{n\geq0}$ is strictly increasing, so the triples obtained with $\varepsilon=\delta=1$ are pairwise distinct.
\end{proof}

For example, the first three members of this subfamily are
\[
 (-55,62,59),\qquad
 (-55,177602,24151),\qquad
 (-55,31879382,4298819).
\]
The independent sign changes in \eqref{eq:family} reflect the symmetries
\[
 z\longmapsto -z,
 \qquad
 y\longmapsto -3026-y
\]
of the equation after $x$ has been fixed at $-55$.

\section{How the construction was found}\label{sec:discovery}

The main obstacle in \eqref{eq:original} is the mixed role of $x$: it occurs both as the coefficient of $z^2$ and, through $x^2$, as a large coefficient of $y$.  Choosing $x$ negative turns the equation into an indefinite quadratic equation in $(y,z)$, which is precisely the setting in which Pell-type propagation becomes available.

Write
\[
 x=-m,
 \qquad m>0.
\]
Then \eqref{eq:original} becomes
\begin{equation}\label{eq:m-equation}
 y^2+(m^2+1)y-mz^2-1=0.
\end{equation}
For odd $m$, completing the square gives
\begin{equation}\label{eq:general-norm}
 \left(y+\frac{m^2+1}{2}\right)^2-mz^2
 =\left(\frac{m^2+1}{2}\right)^2+1.
\end{equation}
Thus, once one finds a single integral point for a suitable nonsquare $m$, a solution of the Pell equation $r^2-ms^2=1$ can be used to generate further points on the same norm level.

The question is therefore how to manufacture a convenient seed.  The dominant terms in \eqref{eq:m-equation} suggest looking for $y$ and $z$ close to $m$.  Put
\[
 y=m+a,
 \qquad
 z=m+b.
\]
Substitution into the left-hand side of \eqref{eq:m-equation} yields
\begin{equation}\label{eq:offset-expansion}
 (a+1-2b)m^2+(2a-b^2+1)m+(a^2+a-1).
\end{equation}
The cubic terms cancel automatically.  We can also force the quadratic term to vanish by taking
\[
 a=2b-1.
\]
Then \eqref{eq:offset-expansion} reduces to
\[
 -(b^2-4b+1)m+4b^2-2b-1.
\]
Consequently, a seed is obtained whenever
\begin{equation}\label{eq:m-rational}
 m=\frac{4b^2-2b-1}{b^2-4b+1}
\end{equation}
is a positive integer.

A particularly natural way to remove the divisibility issue in \eqref{eq:m-rational} is to make the denominator equal to $1$.  Since
\[
 b^2-4b+1=(b-2)^2-3,
\]
the equation $b^2-4b+1=1$ gives $(b-2)^2=4$, hence $b=0$ or $b=4$.  The choice $b=0$ produces $m=-1$, whereas $b=4$ gives
\[
 m=4\cdot4^2-2\cdot4-1=55,
 \qquad
 a=2\cdot4-1=7.
\]
Therefore
\[
 x=-55,
 \qquad
 y=55+7=62,
 \qquad
 z=55+4=59,
\]
which is the seed solution appearing in Theorem~\ref{thm:main}.

For $m=55$, equation \eqref{eq:general-norm} becomes
\[
 (y+1513)^2-55z^2=1513^2+1.
\]
The small identity
\[
 89^2-55\cdot12^2=1
\]
then supplies a norm-one multiplier.  If $U=y+1513$ and $V=z$, multiplication by $89+12\sqrt{55}$ sends
\[
 U+V\sqrt{55}
 \quad\text{to}\quad
 (89U+660V)+(12U+89V)\sqrt{55},
\]
which is exactly the recurrence \eqref{eq:recurrence-step}.  In this way, the seed point and the norm-one identity together explain every coefficient in the final infinite family.

\begin{remark}
Theorem~\ref{thm:main} proves infinitude, as required; it does not assert that the displayed family exhausts all integer solutions of \eqref{eq:original}.
\end{remark}

\end{document}
